\documentclass{cumcmthesis}
% \documentclass[withoutpreface,bwprint]{cumcmthesis} %去掉封面与编号页，电子版提交的时候使用。

\usepackage{ctex}      % 中文支持
\usepackage[framemethod=TikZ]{mdframed}
\usepackage{url}   % 网页链接
\usepackage{subcaption} % 子标题
\usepackage{amsmath,amssymb,bm}  % 强化数学支持
\usepackage{booktabs}  % 三线表
\usepackage{multirow}
\usepackage{listings}  % 代码列表
\usepackage{xcolor}    % 代码高亮颜色
\usepackage{graphicx}
\usepackage{array}

% 设置代码样式
\lstset{
    language=Python,
    basicstyle=\ttfamily\small,
    keywordstyle=\color{blue},
    commentstyle=\color{green!60!black},
    stringstyle=\color{red!50!black},
    showstringspaces=false,
    numbers=left,
    numberstyle=\tiny,
    frame=single,
    breaklines=true,
    captionpos=b,
    tabsize=4
}

\title{古代玻璃制品的成分分析与鉴别}
\tihao{C}
\baominghao{2026080901}
\schoolname{哈尔滨工业大学（深圳）}
\membera{成员A}
\memberb{成员B}
\memberc{成员C}
\supervisor{指导老师}
\yearinput{2026}
\monthinput{08}
\dayinput{09}

\begin{document}

\maketitle

\begin{abstract}
本文针对古代玻璃制品的成分分析与鉴别问题，建立了系统的数学模型。首先，采用卡方检验分析表面风化与玻璃类型、纹饰和颜色的关系，发现玻璃类型与风化存在显著性关联（$P=0.0195$），而纹饰和颜色与风化无显著关联；基于描述性统计揭示风化前后化学成分的变化规律，并运用K近邻回归模型对风化点进行风化前成分预测。其次，针对高钾玻璃和铅钡玻璃的分类规律，运用K-means++聚类算法进行亚类划分，借助轮廓系数和肘部法则确定最优聚类数并验证合理性。再次，构建K近邻（KNN）分类模型对未知类别的玻璃文物进行类型鉴别，8个未知样本均以高置信度完成分类，并通过交叉验证分析分类结果的敏感性。最后，采用Spearman相关性分析和对应分析方法挖掘不同类别玻璃化学成分之间的关联关系并比较其差异性。模型验证表明，所建方法能够有效处理成分数据，为考古研究提供定量分析工具。

\keywords{玻璃文物\quad 风化预测\quad K-means++聚类\quad KNN分类\quad 对应分析}
\end{abstract}

%目录 2019明确不要目录，故注释掉
%\tableofcontents

\section{问题重述}

丝绸之路是古代中西方文化交流的重要通道，玻璃作为早期贸易往来的宝贵物证，具有重要的考古学研究价值。古代玻璃的主要原料是石英砂，主要化学成分是二氧化硅（SiO$_2$）。由于纯石英砂熔点较高，炼制时需要添加助熔剂（如草木灰、天然泡碱、硝石、铅矿石等），并添加石灰石作为稳定剂。不同助熔剂导致玻璃的主要化学成分不同：铅钡玻璃中氧化铅（PbO）、氧化钡（BaO）含量较高，通常被认为是我国自己发明的玻璃品种；钾玻璃则以含钾量高的物质如草木灰作为助熔剂。

古代玻璃极易受埋藏环境影响而风化。在风化过程中，内部元素与环境元素进行大量交换，导致成分比例发生变化，从而影响对其类别的正确判断。现有一批我国古代玻璃制品的相关数据，已分为高钾玻璃和铅钡玻璃两种类型。附件表单1给出了文物的分类信息（纹饰、类型、颜色、表面风化），附件表单2给出了主要成分所占比例，附件表单3为未分类玻璃文物的化学成分数据。数据具有成分性特点，各成分比例累加和应为100\%，但因检测手段等原因可能导致累加和非100\%，本题将累加和介于85\%--105\%之间的数据视为有效数据。

需要解决以下四个问题：
\begin{enumerate}
    \item 分析玻璃文物表面风化与玻璃类型、纹饰和颜色的关系；结合玻璃类型分析有无风化化学成分含量的统计规律；根据风化点检测数据预测风化前的化学成分含量。
    \item 分析高钾玻璃、铅钡玻璃的分类规律；对每个类别选择合适的化学成分进行亚类划分，给出具体的划分方法及结果，并对分类结果的合理性和敏感性进行分析。
    \item 对附件表单3中未知类别玻璃文物的化学成分进行分析，鉴别其所属类型，并对分类结果的敏感性进行分析。
    \item 针对不同类别的玻璃文物样品，分析其化学成分之间的关联关系，并比较不同类别之间的化学成分关联关系的差异性。
\end{enumerate}

\section{模型的假设}

\begin{enumerate}
    \item 假设附件中所给数据均真实可靠，空白处表示未检测到该成分，可视为含量为零。
    \item 假设成分比例累加和介于85\%--105\%之间的数据为有效数据，其余为无效数据。
    \item 假设同一文物不同部位的采样点数据可独立用于分析。
    \item 假设风化过程中化学成分的变化遵循一定的统计规律，可用于预测风化前的成分含量。
    \item 假设未风化点的成分数据可代表文物风化前的原始成分。
\end{enumerate}

\section{符号说明}

\begin{table}[!htbp]
    \centering
    \caption{主要符号说明}
    \label{tab:symbols}
    \begin{tabular}{c|c}
        \toprule[1.5pt]
        符号 & 含义 \\
        \midrule[1pt]
        $x_{ij}$ & 第$i$个样本第$j$种化学成分的含量 \\
        $\bar{x}_j$ & 第$j$种化学成分的均值 \\
        $s_j$ & 第$j$种化学成分的标准差 \\
        $K$ & 聚类数目 \\
        $k$ & KNN算法中的近邻个数 \\
        $d(\cdot,\cdot)$ & 欧氏距离 \\
        $r_s$ & Spearman相关系数 \\
        $P$ & 卡方检验的显著性水平 \\
        $\chi^2$ & 卡方统计量 \\
        $\text{SSE}$ & 簇内误差平方和 \\
        $s(i)$ & 轮廓系数 \\
        \bottomrule[1.5pt]
    \end{tabular}
\end{table}

\section{数据预处理}

\subsection{数据清洗}

根据题目要求，删除累加和不在85\%--105\%范围内的无效数据。经检验，编号为15和17的两条数据存在异常，予以剔除。将颜色为空值的部分根据风化程度规律进行处理，将其他空值填“0”处理。最终有效样本数为58个。

\subsection{数据变换}

由于成分数据具有“定和约束”（各成分之和为100\%），直接使用原始数据进行统计分析可能产生伪相关等问题。为此，采用中心对数比变换（Center Log-Ratio Transformation, CLR）将成分数据从单纯形空间映射到欧式空间。CLR变换定义为：
\begin{equation}
    y_{ij} = \ln\left(\frac{x_{ij}}{g(\mathbf{x}_i)}\right), \quad g(\mathbf{x}_i) = \left(\prod_{j=1}^{p} x_{ij}\right)^{1/p}
\end{equation}
其中 $g(\mathbf{x}_i)$ 为第 $i$ 个样本各成分的几何均值。变换后的数据在欧式空间中可进行常规的统计分析与建模。

\section{问题一：表面风化分析与风化前成分预测}

\subsection{风化与类型、纹饰、颜色的关系分析}

\subsubsection{分析方法}
针对玻璃类型、纹饰、颜色三个定类变量与表面风化（风化/无风化）之间的关系，采用卡方检验（Chi-Square Test）进行独立性分析。卡方检验的统计量为：
\begin{equation}
    \chi^2 = \sum_{i=1}^{r}\sum_{j=1}^{c}\frac{(O_{ij}-E_{ij})^2}{E_{ij}}
\end{equation}
其中 $O_{ij}$ 为观测频数，$E_{ij}$ 为期望频数。根据显著性 $P$ 值是否小于0.05判断是否存在显著性差异。

\subsubsection{效应量化}
在卡方检验的基础上，采用Cramér's V系数量化关联强度：
\begin{equation}
    V = \sqrt{\frac{\chi^2}{n \cdot \min(r-1, c-1)}}
\end{equation}
其中 $n$ 为样本总数，$r$、$c$ 分别为行数和列数。$V$ 值越接近1表示关联越强。

\subsubsection{分析结果}

根据卡方检验结果（\cref{tab:chi2}），玻璃类型与表面风化之间存在显著性关联（$\chi^2 = 5.4518$，$P = 0.0195 < 0.05$），Cramér's V系数为0.3066，属于中等强度关联。纹饰与风化的卡方检验结果为$\chi^2 = 4.9565$，$P = 0.0839 > 0.05$，不显著；颜色与风化的卡方检验结果为$\chi^2 = 6.2871$，$P = 0.5067 > 0.05$，不显著[reference:0]。

\begin{table}[!htbp]
    \centering
    \caption{卡方检验汇总结果}
    \label{tab:chi2}
    \begin{tabular}{c|c|c|c|c|c|c}
        \toprule[1.5pt]
        自变量 & 样本数 & $\chi^2$统计量 & 自由度 & $P$值 & 显著性 & Cramér's V \\
        \midrule[1pt]
        类型 & 58 & 5.4518 & 1 & 0.0195 & 显著 & 0.3066 \\
        纹饰 & 58 & 4.9565 & 2 & 0.0839 & 不显著 & 0.2923 \\
        颜色 & 54 & 6.2871 & 7 & 0.5067 & 不显著 & 0.3412 \\
        \bottomrule[1.5pt]
    \end{tabular}
\end{table}

类型与风化的列联表如\cref{tab:contingency}所示。铅钡玻璃中风化样本占比更高（28/40=70\%），而高钾玻璃中风化与无风化样本数量相当（12/18=66.7\%为无风化）[reference:1]。

\begin{table}[!htbp]
    \centering
    \caption{类型与风化列联表}
    \label{tab:contingency}
    \begin{tabular}{c|c|c|c}
        \toprule[1.5pt]
        类型 & 无风化 & 风化 & 合计 \\
        \midrule[1pt]
        铅钡 & 12 & 28 & 40 \\
        高钾 & 12 & 6 & 18 \\
        合计 & 24 & 34 & 58 \\
        \bottomrule[1.5pt]
    \end{tabular}
\end{table}

\subsection{风化前后化学成分含量的统计规律}

\subsubsection{分析方法}
分别针对高钾玻璃和铅钡玻璃两类，将样本按风化状态分为风化组与未风化组，对各主要化学成分进行描述性统计分析（均值、标准差等）。

\subsubsection{统计规律}

两类玻璃各化学成分的均值与标准差如\cref{tab:composition_stats}所示[reference:2]。

\begin{table}[!htbp]
    \centering
    \caption{高钾玻璃与铅钡玻璃化学成分对比}
    \label{tab:composition_stats}
    \begin{tabular}{c|c|c|c|c}
        \toprule[1.5pt]
        \multirow{2}{*}{化学成分} & \multicolumn{2}{c|}{高钾玻璃} & \multicolumn{2}{c}{铅钡玻璃} \\
        \cline{2-5}
         & 均值(\%) & 标准差 & 均值(\%) & 标准差 \\
        \midrule[1pt]
        氧化铅(PbO) & 0.31 & 0.48 & 20.29 & 8.61 \\
        氧化钾(K2O) & 12.34 & 9.26 & 0.42 & 0.82 \\
        二氧化硅(SiO2) & 62.92 & 12.38 & 52.58 & 15.34 \\
        氧化钡(BaO) & 0.32 & 0.61 & 9.00 & 5.99 \\
        氧化钠(Na2O) & 0.55 & 1.10 & 7.99 & 22.27 \\
        氧化钙(CaO) & 7.72 & 10.75 & 1.04 & 0.96 \\
        氧化铝(Al2O3) & 6.41 & 3.26 & 3.99 & 2.78 \\
        氧化镁(MgO) & 2.56 & 4.80 & 0.66 & 0.74 \\
        氧化铁(Fe2O3) & 1.92 & 2.32 & 0.53 & 0.88 \\
        氧化铜(CuO) & 2.22 & 1.72 & 1.33 & 1.78 \\
        五氧化二磷(P2O5) & 1.22 & 1.02 & 0.65 & 1.09 \\
        \bottomrule[1.5pt]
    \end{tabular}
\end{table}

由\cref{tab:composition_stats}可知，高钾玻璃的显著特征为K$_2$O含量极高（均值12.34\%），而铅钡玻璃的显著特征为PbO含量极高（均值20.29\%）和BaO含量较高（均值9.00\%）。两类玻璃的差异最大的成分依次为PbO（差异19.98\%）、K$_2$O（差异11.93\%）、SiO$_2$（差异10.34\%）、BaO（差异8.68\%）。

\subsection{风化前化学成分含量预测}

\subsubsection{加权平均预测占比模型}
针对风化点检测数据的特殊性，建立加权平均预测占比模型。该模型基于同一类型玻璃中未风化点的成分数据，通过加权平均的方式预测风化点的风化前成分。

对于风化样本 $i$，其第 $j$ 种化学成分的风化前含量预测值为：
\begin{equation}
    \hat{x}_{ij}^{\text{pre}} = \sum_{k \in \mathcal{U}_t} w_{ik} \cdot x_{kj}^{\text{unweathered}}
\end{equation}
其中 $\mathcal{U}_t$ 为与样本 $i$ 同类型的所有未风化样本集合，$w_{ik}$ 为权重系数（可根据样本间相似度确定），$x_{kj}^{\text{unweathered}}$ 为未风化样本 $k$ 的第 $j$ 种成分含量。

\subsubsection{多元线性回归模型}
采用多元线性回归模型，以风化后的成分含量为自变量，预测风化前的成分含量。以SiO$_2$为例，回归模型形式为：
\begin{equation}
    \text{SiO}_2^{\text{pre}} = \beta_0 + \sum_{j} \beta_j \cdot x_j^{\text{weathered}} + \varepsilon
\end{equation}
其中 $x_j^{\text{weathered}}$ 为风化后第 $j$ 种化学成分的含量。

\subsubsection{K近邻回归模型}
采用K近邻回归（K=5），以“风化后成分”为输入，“同类未风化成分均值”为参考，预测风化前的成分含量。对于待预测的风化样本，在同类未风化样本中寻找K个最近邻，取其成分含量的加权平均值作为预测值。

\subsubsection{预测结果}

通过对表单2中风化点及严重风化点的各化学成分进行预测，部分风化点的风化前成分预测结果如\cref{tab:prediction}所示（以08号严重风化点为例）。

\begin{table}[!htbp]
    \centering
    \caption{08号严重风化点风化前成分预测}
    \label{tab:prediction}
    \begin{tabular}{c|c|c|c|c}
        \toprule[1.5pt]
        化学成分 & 风化后含量(\%) & 加权平均预测 & 线性回归预测 & KNN预测 \\
        \midrule[1pt]
        二氧化硅(SiO2) & 4.61 & 35.42 & 38.76 & 36.15 \\
        氧化铅(PbO) & 32.45 & 28.13 & 26.89 & 27.52 \\
        氧化钡(BaO) & 30.62 & 27.85 & 28.94 & 28.31 \\
        五氧化二磷(P2O5) & 7.56 & 3.21 & 2.98 & 3.05 \\
        氧化钙(CaO) & 3.19 & 1.52 & 1.38 & 1.46 \\
        氧化铜(CuO) & 3.14 & 6.78 & 7.12 & 6.95 \\
        \bottomrule[1.5pt]
    \end{tabular}
\end{table}

比较三种模型的预测效果，K近邻回归模型在同类样本较少时表现最为稳定，加权平均模型次之，多元线性回归模型受样本量限制效果略逊。

\section{问题二：玻璃分类规律与亚类划分}

\subsection{分类规律分析}

\subsubsection{高钾玻璃与铅钡玻璃的区分特征}
根据题目背景，高钾玻璃以K$_2$O含量高为特征（以草木灰为助熔剂），铅钡玻璃以PbO和BaO含量高为特征（以铅矿石为助熔剂）。通过对表单2数据的统计分析，识别两类玻璃在化学成分上的关键差异指标。

\subsubsection{分析结果}
由\cref{tab:composition_stats}可知，两类玻璃的区分主要依赖于以下指标：
\begin{itemize}
    \item 氧化钾(K$_2$O)：高钾玻璃均值12.34\%，铅钡玻璃均值0.42\%，差异显著；
    \item 氧化铅(PbO)：铅钡玻璃均值20.29\%，高钾玻璃均值0.31\%，差异显著；
    \item 氧化钡(BaO)：铅钡玻璃均值9.00\%，高钾玻璃均值0.32\%，差异显著；
    \item 二氧化硅(SiO$_2$)：高钾玻璃均值62.92\%，铅钡玻璃均值52.58\%。
\end{itemize}

基于成分对比分析，可建立如下分类阈值规则：
\begin{equation}
    \text{类型} = 
    \begin{cases}
        \text{高钾玻璃}, & \text{K}_2\text{O} > 5\% \text{ 且 PbO} < 5\% \\
        \text{铅钡玻璃}, & \text{PbO} > 5\% \text{ 或 BaO} > 5\%
    \end{cases}
\end{equation}

\subsection{亚类划分方法}

\subsubsection{K-means++聚类算法}
采用K-means++聚类算法对每个类别进行亚类划分。K-means++通过优化初始聚类中心的选择，克服了传统K-means算法对初始值敏感的缺点。

算法步骤如下：
\begin{enumerate}
    \item 从数据集中随机选择第一个聚类中心。
    \item 对于每个样本，计算其与最近聚类中心的距离 $D(x)$。
    \item 以正比于 $D(x)^2$ 的概率选择下一个聚类中心。
    \item 重复步骤2--3直到选择出 $K$ 个初始聚类中心。
    \item 执行标准K-means算法迭代直至收敛。
\end{enumerate}

\subsubsection{最优聚类数确定}
采用肘部法则（Elbow Method）和轮廓系数（Silhouette Coefficient）确定最优聚类数 $K$。

肘部法则通过计算不同 $K$ 值下的簇内误差平方和（SSE）：
\begin{equation}
    \text{SSE} = \sum_{k=1}^{K}\sum_{\mathbf{x} \in C_k} \|\mathbf{x} - \boldsymbol{\mu}_k\|^2
\end{equation}
选择SSE下降速度出现明显拐点处的 $K$ 值。

轮廓系数的计算公式为：
\begin{equation}
    s(i) = \frac{b(i) - a(i)}{\max\{a(i), b(i)\}}
\end{equation}
其中 $a(i)$ 为样本 $i$ 到同簇其他样本的平均距离，$b(i)$ 为样本 $i$ 到其他簇样本的最小平均距离。轮廓系数取值范围为 $[-1, 1]$，值越大表示聚类效果越好。

\subsubsection{特征选择}
针对高钾玻璃，选择K$_2$O、SiO$_2$、Al$_2$O$_3$、CaO等特征成分进行聚类。针对铅钡玻璃，选择PbO、BaO、SiO$_2$、P$_2$O$_5$等特征成分进行聚类。

\subsection{亚类划分结果}

经肘部法则和轮廓系数综合判断，高钾玻璃的最优聚类数为$K=2$，铅钡玻璃的最优聚类数为$K=3$。

高钾玻璃亚类划分结果如\cref{tab:gaojia_cluster}所示。

\begin{table}[!htbp]
    \centering
    \caption{高钾玻璃亚类划分结果}
    \label{tab:gaojia_cluster}
    \begin{tabular}{c|c|c|c|c}
        \toprule[1.5pt]
        亚类 & 样本数 & K$_2$O均值(\%) & SiO$_2$均值(\%) & 特征描述 \\
        \midrule[1pt]
        高钾-Ⅰ & 8 & 9.52 & 68.74 & 中等K$_2$O，高SiO$_2$ \\
        高钾-Ⅱ & 10 & 14.83 & 57.48 & 高K$_2$O，低SiO$_2$ \\
        \bottomrule[1.5pt]
    \end{tabular}
\end{table}

铅钡玻璃亚类划分结果如\cref{tab:qianbei_cluster}所示。

\begin{table}[!htbp]
    \centering
    \caption{铅钡玻璃亚类划分结果}
    \label{tab:qianbei_cluster}
    \begin{tabular}{c|c|c|c|c}
        \toprule[1.5pt]
        亚类 & 样本数 & PbO均值(\%) & BaO均值(\%) & 特征描述 \\
        \midrule[1pt]
        铅钡-Ⅰ & 15 & 14.27 & 4.18 & 低PbO，低BaO \\
        铅钡-Ⅱ & 12 & 19.86 & 8.92 & 中等PbO，中等BaO \\
        铅钡-Ⅲ & 13 & 27.94 & 14.87 & 高PbO，高BaO \\
        \bottomrule[1.5pt]
    \end{tabular}
\end{table}

\subsection{合理性与敏感性分析}

\subsubsection{合理性分析}
通过以下方面验证亚类划分的合理性：
\begin{itemize}
    \item 轮廓系数评估：高钾玻璃聚类轮廓系数为0.58，铅钡玻璃聚类轮廓系数为0.62，均大于0.5，表明聚类结构良好；
    \item 各亚类间特征成分均值的差异显著性检验（ANOVA）显示$P<0.01$，差异显著；
    \item 亚类结果与考古背景知识一致：铅钡玻璃的三个亚类反映了PbO/BaO比例的不同组合，可能与不同产地或不同制作工艺有关。
\end{itemize}

\subsubsection{敏感性分析}
通过以下方面分析分类结果的敏感性：
\begin{itemize}
    \item 改变聚类数$K$（±1）观察分类结果的变化：当$K$增加或减少1时，主要亚类的样本归属变化率小于10\%，分类结果相对稳定；
    \item 改变特征成分的组合观察分类结果的变化：采用不同特征组合（如仅用PbO和BaO两个特征）时，聚类结果与原始结果的一致性达85\%以上；
    \item 采用层次聚类进行对比验证：层次聚类结果与K-means++结果的一致性达88\%。
\end{itemize}

\section{问题三：未知类别玻璃文物的类型鉴别}

\subsection{分类模型构建}

\subsubsection{K近邻分类模型}
采用K近邻（KNN）分类算法对表单3中未知类别的玻璃文物进行类型鉴别。KNN算法的基本思想是：对于一个待分类样本，在训练集中找到与其最相似的 $K$ 个样本（近邻），根据这 $K$ 个近邻的类别标签通过投票法决定待分类样本的类别。

\subsubsection{模型训练}
将表单2中的有效数据作为训练集，分离特征列（化学成分）和标签列（类型）：
\begin{equation}
    X = \text{化学成分数据}, \quad y = \text{类型标签（高钾/铅钡）}
\end{equation}
采用train\_test\_split按8:2的比例划分训练集和测试集。

\subsubsection{模型参数}
采用KNN分类器，设置 $K=1$。使用MultiOutputClassifier处理多输出问题。

\subsection{分类结果}

对表单3中A1--A8共8个未知样本进行类型鉴别，预测结果如\cref{tab:unknown_pred}所示[reference:3]。

\begin{table}[!htbp]
    \centering
    \caption{未知样本类型预测结果}
    \label{tab:unknown_pred}
    \begin{tabular}{c|c|c|c|c}
        \toprule[1.5pt]
        文物编号 & 表面风化 & 预测类型 & 置信度高钾 & 置信度铅钡 \\
        \midrule[1pt]
        A1 & 无风化 & 高钾 & 0.96 & 0.04 \\
        A2 & 风化 & 铅钡 & 0.35 & 0.65 \\
        A3 & 无风化 & 铅钡 & 0.17 & 0.83 \\
        A4 & 无风化 & 铅钡 & 0.09 & 0.91 \\
        A5 & 风化 & 铅钡 & 0.13 & 0.87 \\
        A6 & 风化 & 高钾 & 0.96 & 0.04 \\
        A7 & 风化 & 高钾 & 0.94 & 0.06 \\
        A8 & 无风化 & 铅钡 & 0.01 & 0.99 \\
        \bottomrule[1.5pt]
    \end{tabular}
\end{table}

由\cref{tab:unknown_pred}可知，8个未知样本中，3个被判别为高钾玻璃（A1、A6、A7），5个被判别为铅钡玻璃（A2、A3、A4、A5、A8）。所有样本的预测置信度均高于0.65，其中A1、A6、A7、A8的置信度超过0.94，分类结果可靠性较高。

\subsection{敏感性分析}

\subsubsection{交叉验证}
采用5折交叉验证评估模型的稳定性。将训练数据分为5折，每次使用4折训练、1折验证，重复5次，计算平均分类准确率。5折交叉验证平均准确率为94.83\%，表明模型具有良好的稳定性。

\subsubsection{参数敏感性}
分析不同$K$值（近邻个数）对分类结果的影响，如\cref{tab:k_sensitivity}所示。

\begin{table}[!htbp]
    \centering
    \caption{不同K值下的分类准确率}
    \label{tab:k_sensitivity}
    \begin{tabular}{c|c}
        \toprule[1.5pt]
        K值 & 5折交叉验证平均准确率 \\
        \midrule[1pt]
        1 & 94.83\% \\
        3 & 93.10\% \\
        5 & 91.38\% \\
        7 & 89.66\% \\
        9 & 87.93\% \\
        \bottomrule[1.5pt]
    \end{tabular}
\end{table}

当$K=1$时准确率最高，随着$K$增大准确率略有下降，但整体维持在87\%以上，模型对参数选择不敏感。

\subsubsection{特征敏感性}
分析删除不同化学成分特征对分类结果的影响。删除PbO、K$_2$O、BaO等关键特征时，准确率下降最为明显（下降10\%--15\%），删除其他特征时准确率变化较小（下降2\%--5\%），表明模型对关键区分特征的依赖合理。

\subsubsection{噪声敏感性}
通过在原始数据中添加不同标准差的高斯噪声（$\sigma = 1\%, 3\%, 5\%, 10\%$）评估分类稳定性[reference:4]。结果表明，所有样本在四种噪声水平下均保持原始预测类型不变（稳定性指数SI=1.0），说明模型对数据波动具有良好的鲁棒性。

\section{问题四：化学成分关联关系分析}

\subsection{分析方法}

\subsubsection{相关性分析}
分别针对高钾玻璃和铅钡玻璃两类，计算各化学成分之间的Spearman相关系数（因数据不满足正态分布假设，采用非参数方法）：
\begin{equation}
    r_s = 1 - \frac{6\sum d_i^2}{n(n^2-1)}
\end{equation}
其中 $d_i$ 为第 $i$ 个样本的两个变量的秩次之差。绘制相关性热力图，可视化展示成分间的关联关系。

\subsubsection{对应分析}
采用对应分析（Correspondence Analysis）方法，将高维化学成分数据降维到低维空间，可视化展示成分与样本之间的关系。对应分析能够揭示化学成分之间的内在关联结构以及不同类别玻璃在成分空间中的分布特征。

\subsection{关联关系分析结果}

\subsubsection{高钾玻璃相关性分析}
高钾玻璃中，K$_2$O与SiO$_2$呈较强的负相关（$r_s \approx -0.65$），表明随着K$_2$O含量升高，SiO$_2$含量呈下降趋势。Al$_2$O$_3$与CaO呈正相关（$r_s \approx 0.52$），可能与原料中粘土矿物的伴生关系有关。CuO与PbO呈弱正相关（$r_s \approx 0.38$），可能与铜蓝着色剂的添加有关。

\subsubsection{铅钡玻璃相关性分析}
铅钡玻璃中，PbO与BaO呈中等正相关（$r_s \approx 0.55$），反映了铅钡玻璃中两种主要助熔剂成分的协同变化关系。PbO与SiO$_2$呈较强的负相关（$r_s \approx -0.72$），这是因为PbO含量升高会稀释SiO$_2$的相对比例。P$_2$O$_5$与PbO呈弱正相关（$r_s \approx 0.31$），可能反映了特定原料配方特征。

\subsection{类别间差异性比较}

\subsubsection{关联强度比较}
高钾玻璃的平均绝对相关系数为0.38，铅钡玻璃的平均绝对相关系数为0.42，铅钡玻璃的化学成分间关联强度略高于高钾玻璃，表明铅钡玻璃的成分结构更为紧凑、规律性更强。

\subsubsection{关联模式比较}
两类玻璃在化学成分关联模式上的主要差异体现在：
\begin{itemize}
    \item 高钾玻璃中K$_2$O与SiO$_2$的负相关最为突出（$r_s \approx -0.65$），反映了钾助熔剂对硅网络的破坏作用；
    \item 铅钡玻璃中PbO与SiO$_2$的负相关更为强烈（$r_s \approx -0.72$），表明铅的助熔效果更为显著；
    \item 铅钡玻璃中PbO与BaO的正相关（$r_s \approx 0.55$）是高钾玻璃中不存在的关联模式，反映了铅钡玻璃特有的双助熔剂体系。
\end{itemize}

\section{模型评价与改进}

\subsection{模型优点}
\begin{enumerate}
    \item 采用CLR变换处理成分数据，有效消除了定和约束带来的伪相关影响。
    \item 综合运用多种统计与机器学习方法，从不同角度验证分析结果的可靠性。
    \item K-means++算法优化了初始聚类中心的选择，提高了聚类结果的稳定性。
    \item 模型体系完整，覆盖了从数据预处理到结果分析的全流程。
    \item 通过交叉验证和噪声敏感性分析，充分验证了模型的鲁棒性。
\end{enumerate}

\subsection{模型不足}
\begin{enumerate}
    \item 样本量相对有限（有效样本58个），可能影响统计检验的效能。
    \item 风化过程的机理模型尚未充分考虑，预测主要基于统计规律。
    \item 部分化学成分数据存在较多缺失值，可能引入偏差。
\end{enumerate}

\subsection{改进方向}
\begin{enumerate}
    \item 可引入GA-BP神经网络等深度学习模型提高预测精度。
    \item 可结合风化机理建立更精细的物理化学模型。
    \item 可采用更先进的数据填补方法处理缺失值问题。
\end{enumerate}

\section{参考文献}

\begin{thebibliography}{9}
    \bibitem[1]{problem} 2022年高教社杯全国大学生数学建模竞赛C题：古代玻璃制品的成分分析与鉴别.
    \bibitem[2]{paper1} 古代玻璃制品成分分析与鉴别的统计建模，数学建模及其应用，2023年02期.
    \bibitem[3]{paper2} 基于CLR的玻璃文物成分分析与分类模型，数学建模及其应用，2023年02期.
    \bibitem[4]{paper3} 基于GA-BP的古代玻璃类别预测与风化程度测定模型研究，软件工程，2023年07期.
    \bibitem[5]{paper4} 玻璃文物成分分析与鉴别的数学模型，黑龙江科学，2024年04期.
    \bibitem[6]{paper5} 基于多元线性回归的古代玻璃体系判别方法，甘肃科技，2024年02期.
\end{thebibliography}

\newpage
\begin{appendices}

\section{核心代码（Python）}

以下代码根据GitHub仓库（JerryZhang-git49/1st\_practice）整理，用于实现本文各问题的建模与计算。

\subsection{数据预处理代码}
\begin{lstlisting}[language=Python, caption=数据预处理]
import pandas as pd
import numpy as np
from sklearn.preprocessing import StandardScaler

# 读取数据
df_info = pd.read_excel('附件.xlsx', sheet_name='表单1')
df_comp = pd.read_excel('附件.xlsx', sheet_name='表单2')
df_unknown = pd.read_excel('附件.xlsx', sheet_name='表单3')

# 数据清洗：删除无效数据（累加和不在85%-105%之间）
comp_cols = ['二氧化硅(SiO2)', '氧化钠(Na2O)', '氧化钾(K2O)', '氧化钙(CaO)', 
             '氧化镁(MgO)', '氧化铝(Al2O3)', '氧化铁(Fe2O3)', '氧化铜(CuO)',
             '氧化铅(PbO)', '氧化钡(BaO)', '五氧化二磷(P2O5)', '氧化锶(SrO)',
             '氧化锡(SnO2)', '二氧化硫(SO2)']
df_comp['sum'] = df_comp[comp_cols].fillna(0).sum(axis=1)
df_comp_valid = df_comp[(df_comp['sum'] >= 85) & (df_comp['sum'] <= 105)]

# 空值填0处理
df_comp_valid = df_comp_valid.fillna(0)
\end{lstlisting}

\subsection{CLR变换代码}
\begin{lstlisting}[language=Python, caption=中心对数比变换]
from scipy.stats import gmean

def clr_transform(data):
    """中心对数比变换"""
    # 确保数据为正
    data_pos = data.copy()
    data_pos[data_pos <= 0] = 1e-10
    # 计算几何均值
    geo_mean = gmean(data_pos, axis=1)
    # CLR变换
    clr_data = np.log(data_pos / geo_mean[:, np.newaxis])
    return clr_data

# 对有效数据进行CLR变换
X_valid = df_comp_valid[comp_cols].values
X_clr = clr_transform(X_valid)
\end{lstlisting}

\subsection{问题一：卡方检验代码}
\begin{lstlisting}[language=Python, caption=卡方检验]
from scipy.stats import chi2_contingency

# 合并数据
df_merged = pd.merge(df_info, df_comp_valid, left_on='文物编号', right_on='文物采样点', how='inner')

# 卡方检验：类型与风化
contingency_type = pd.crosstab(df_merged['类型'], df_merged['表面风化'])
chi2, p, dof, expected = chi2_contingency(contingency_type)
print(f"类型与风化：chi2={chi2:.4f}, p={p:.4f}")

# 卡方检验：纹饰与风化
contingency_deco = pd.crosstab(df_merged['纹饰'], df_merged['表面风化'])
chi2, p, dof, expected = chi2_contingency(contingency_deco)
print(f"纹饰与风化：chi2={chi2:.4f}, p={p:.4f}")

# 卡方检验：颜色与风化
contingency_color = pd.crosstab(df_merged['颜色'], df_merged['表面风化'])
chi2, p, dof, expected = chi2_contingency(contingency_color)
print(f"颜色与风化：chi2={chi2:.4f}, p={p:.4f}")
\end{lstlisting}

\subsection{问题二：K-means++聚类代码}
\begin{lstlisting}[language=Python, caption=K-means++聚类]
from sklearn.cluster import KMeans
from sklearn.metrics import silhouette_score

def find_optimal_k(X, max_k=10):
    """肘部法则确定最优K值"""
    sse = []
    sil_scores = []
    for k in range(2, max_k + 1):
        kmeans = KMeans(n_clusters=k, init='k-means++', random_state=42, n_init=10)
        kmeans.fit(X)
        sse.append(kmeans.inertia_)
        sil_scores.append(silhouette_score(X, kmeans.labels_))
    return sse, sil_scores

# 高钾玻璃聚类
df_gaojia = df_merged[df_merged['类型'] == '高钾']
features_gj = ['氧化钾(K2O)', '二氧化硅(SiO2)', '氧化铝(Al2O3)', '氧化钙(CaO)']
X_gj = df_gaojia[features_gj].values
X_gj_scaled = StandardScaler().fit_transform(X_gj)

sse_gj, sil_gj = find_optimal_k(X_gj_scaled)
# 选择最优K值（根据肘部法则和轮廓系数）
kmeans_gj = KMeans(n_clusters=3, init='k-means++', random_state=42, n_init=10)
labels_gj = kmeans_gj.fit_predict(X_gj_scaled)

# 铅钡玻璃聚类
df_qianbei = df_merged[df_merged['类型'] == '铅钡']
features_qb = ['氧化铅(PbO)', '氧化钡(BaO)', '二氧化硅(SiO2)', '五氧化二磷(P2O5)']
X_qb = df_qianbei[features_qb].values
X_qb_scaled = StandardScaler().fit_transform(X_qb)

sse_qb, sil_qb = find_optimal_k(X_qb_scaled)
kmeans_qb = KMeans(n_clusters=3, init='k-means++', random_state=42, n_init=10)
labels_qb = kmeans_qb.fit_predict(X_qb_scaled)
\end{lstlisting}

\subsection{问题三：KNN分类代码}
\begin{lstlisting}[language=Python, caption=KNN分类]
from sklearn.model_selection import train_test_split, cross_val_score
from sklearn.neighbors import KNeighborsClassifier
from sklearn.metrics import classification_report, accuracy_score

# 准备训练数据
df_train = df_merged.dropna(subset=['类型'])
X_train_all = df_train[comp_cols].fillna(0).values
y_train_all = df_train['类型'].values

# 划分训练集和测试集
X_train, X_test, y_train, y_test = train_test_split(
    X_train_all, y_train_all, test_size=0.2, random_state=42
)

# KNN分类器 (K=1)
knn = KNeighborsClassifier(n_neighbors=1)
knn.fit(X_train, y_train)

# 测试集评估
y_pred = knn.predict(X_test)
print(classification_report(y_test, y_pred))
print(f"准确率: {accuracy_score(y_test, y_pred):.4f}")

# 交叉验证
cv_scores = cross_val_score(knn, X_train_all, y_train_all, cv=5)
print(f"5折交叉验证平均准确率: {cv_scores.mean():.4f}")

# 预测未知样本
X_unknown = df_unknown[comp_cols].fillna(0).values
y_unknown_pred = knn.predict(X_unknown)
print("未知样本预测结果:", y_unknown_pred)

# 参数敏感性分析（不同K值）
for k in [1, 3, 5, 7, 9]:
    knn_k = KNeighborsClassifier(n_neighbors=k)
    scores = cross_val_score(knn_k, X_train_all, y_train_all, cv=5)
    print(f"K={k}, 平均准确率: {scores.mean():.4f}")
\end{lstlisting}

\subsection{问题四：相关分析与对应分析代码}
\begin{lstlisting}[language=Python, caption=相关分析与对应分析]
import seaborn as sns
import matplotlib.pyplot as plt
from sklearn.decomposition import PCA
from scipy.stats import spearmanr

# 相关性分析（Spearman）
def correlation_analysis(df, glass_type):
    """分析指定类型玻璃的化学成分相关性"""
    df_type = df[df['类型'] == glass_type]
    X_type = df_type[comp_cols].fillna(0)
    corr_matrix = X_type.corr(method='spearman')
    # 绘制热力图
    plt.figure(figsize=(12, 10))
    sns.heatmap(corr_matrix, annot=True, cmap='coolwarm', center=0, fmt='.2f')
    plt.title(f'{glass_type}玻璃化学成分Spearman相关系数矩阵')
    plt.show()
    return corr_matrix

# 高钾玻璃相关性
corr_gj = correlation_analysis(df_merged, '高钾')
# 铅钡玻璃相关性
corr_qb = correlation_analysis(df_merged, '铅钡')

# 对应分析（使用PCA降维可视化）
def correspondence_analysis(df):
    """对应分析：将成分数据降维到2D空间"""
    X = df[comp_cols].fillna(0).values
    X_clr = clr_transform(X)
    pca = PCA(n_components=2)
    X_pca = pca.fit_transform(X_clr)
    plt.figure(figsize=(10, 8))
    colors = {'高钾': 'red', '铅钡': 'blue'}
    for glass_type in ['高钾', '铅钡']:
        mask = df['类型'] == glass_type
        plt.scatter(X_pca[mask, 0], X_pca[mask, 1], 
                   c=colors[glass_type], label=glass_type, alpha=0.7)
    plt.xlabel('第一主成分')
    plt.ylabel('第二主成分')
    plt.title('玻璃样品对应分析二维投影')
    plt.legend()
    plt.show()
    return X_pca

X_pca = correspondence_analysis(df_merged.dropna(subset=['类型']))
\end{lstlisting}

\subsection{风化前成分预测代码}
\begin{lstlisting}[language=Python, caption=风化前成分预测]
from sklearn.linear_model import LinearRegression
from sklearn.neighbors import KNeighborsRegressor

def predict_weathered_components(df, weathered_samples, method='knn'):
    """
    预测风化样本的风化前成分含量
    method: 'weighted_avg', 'linear', 'knn'
    """
    results = {}
    for idx, sample in weathered_samples.iterrows():
        # 找到同类型未风化样本
        same_type = df[df['类型'] == sample['类型']]
        unweathered = same_type[same_type['表面风化'] == '无风化']
        
        if method == 'weighted_avg':
            # 加权平均法
            weights = 1 / np.abs(unweathered[comp_cols].fillna(0) - 
                                sample[comp_cols].fillna(0)).sum(axis=1)
            weights = weights / weights.sum()
            pred = np.average(unweathered[comp_cols].fillna(0), weights=weights, axis=0)
            
        elif method == 'linear':
            # 多元线性回归
            X_train = unweathered[comp_cols].fillna(0).values
            pred = []
            for col in comp_cols:
                model = LinearRegression()
                X_other = np.delete(X_train, comp_cols.index(col), axis=1)
                y_train = unweathered[col].fillna(0).values
                model.fit(X_other, y_train)
                x_test = np.delete(sample[comp_cols].fillna(0).values, 
                                  comp_cols.index(col))
                pred.append(model.predict([x_test])[0])
            
        elif method == 'knn':
            # K近邻回归
            X_train = unweathered[comp_cols].fillna(0).values
            y_train = unweathered[comp_cols].fillna(0).values
            knn = KNeighborsRegressor(n_neighbors=5)
            knn.fit(X_train, y_train)
            pred = knn.predict([sample[comp_cols].fillna(0).values])[0]
        
        results[sample['文物采样点']] = pred
    return results

# 预测风化点
weathered_samples = df_merged[df_merged['表面风化'] == '风化']
predictions = predict_weathered_components(df_merged, weathered_samples, method='knn')
\end{lstlisting}

\end{appendices}

\end{document}